\section{Introduction}
\label{sec:introduction}

Autonomous driving systems are increasingly built around the idea that an
agent should be able to \emph{imagine} the consequences of its actions before
committing to them. A \emph{world model} provides exactly this capability: it
learns the dynamics of the environment from observation and can roll out
plausible future states conditioned on the past. When the observations are
front-camera videos, such a model becomes a learned, data-driven simulator of
the road ahead. The appeal is twofold. First, large amounts of raw driving
video are available without any manual annotation, so a world model can in
principle absorb the long tail of real driving behaviour that is expensive to
label. Second, a model that predicts future frames implicitly has to learn
perception, scene dynamics, and the ego-vehicle's own motion, which makes its
internal representation a natural substrate for downstream decision making.

This report studies the internal representation of one such world model,
\textbf{Orbis}~\cite{orbis}, and asks a concrete question: \emph{does the latent
produced by a driving world model already contain the information needed to
predict where the car is going?} To answer it, we attach a lightweight
\emph{attentive diffusion probe} to the frozen image latent of Orbis and train
it to predict the future 2D trajectory of the ego vehicle. The probe is
deliberately small and non-invasive so that any predictive success can be
attributed to the world-model representation rather than to the probe itself.
We then assess both the pointwise accuracy and the \emph{distribution} of the
predicted trajectories, since a good world model should capture the
multi-modal, uncertain nature of driving rather than collapse to a single
averaged path.

The remainder of this report provides the background needed to situate this
study. This introduction describes the Orbis world model and the proposed
probing setup. Section~\ref{sec:related} reviews the theory of diffusion and
flow matching, the diffusion transformer architecture that underlies modern
generative models, and the broader landscape of world models. Section~\ref{sec:datasets}
describes the data and benchmarks --- nuPlan and the NAVSIM pseudo-simulation
benchmark --- on which driving trajectories are produced and scored. The
experimental part of the project is written up separately.

\subsection{The Orbis world model}
\label{sec:orbis}

Orbis~\cite{orbis} is a driving world model that predicts future video frames
from front-camera observations. Its central finding is that a continuous,
flow-matching world model is both more economical and more robust on
long-horizon rollouts than its discrete, token-based counterpart: with only
$469$M parameters and roughly $280$ hours of raw video, Orbis reaches
state-of-the-art video-prediction quality on the nuPlan, Waymo, and
nuPlan-turns benchmarks, while remaining stable on rollouts of up to twenty
seconds. Crucially for our purposes, the model is trained on raw video alone,
without depth maps, bounding boxes, HD maps, or any other privileged
supervision, so its latent representation is learned end-to-end from the
predictive objective.

Architecturally, Orbis follows the now-standard two-stage recipe for latent
generative modelling, illustrated in Figure~\ref{fig:orbis}. \textbf{Stage~1}
is an image \emph{tokenizer} --- an autoencoder in the spirit of a VAE --- that
compresses a frame into a compact latent and decodes it back to pixels.
\textbf{Stage~2} is the latent \emph{world model} that operates entirely in this
latent space, autoregressively generating the next frame's latent given a
window of context frames. Decoupling the two stages lets the expensive
generative model work at a low spatial resolution while the tokenizer handles
the high-frequency pixel detail.

\paragraph{Stage 1: hybrid image tokenizer.}
The tokenizer maps an image $\mathbf{I}\in\mathbb{R}^{H\times W\times 3}$ to a
spatial latent compressed by a factor of $16$ in each dimension. A distinctive
feature of Orbis is its \emph{hybrid, factorized} design. Two separate encoders
produce complementary representations: a \emph{semantic} encoder
$\mathcal{E}_s$, regularized to stay close to DINOv2 features through a cosine
similarity loss, and a \emph{detail} encoder $\mathcal{E}_d$ that captures
low-level appearance. Each stream is quantized against its own codebook,
yielding semantic tokens $\mathbf{q}_s$ and detail tokens $\mathbf{q}_d$, and
the tokenizer is trained with the VQGAN objective. By fine-tuning with a
$50\%$ probability of bypassing the vector-quantization bottleneck, the same
encoder supports \emph{both} a continuous latent
$\mathbf{x}=(\mathbf{x}_s;\mathbf{x}_d)$ and a discrete latent
$\mathbf{q}=(\mathbf{q}_s;\mathbf{q}_d)$. This hybrid construction is what
allows Orbis to compare continuous (flow-matching) and discrete (masked
generative) world models on an identical footing. In this report the
\emph{continuous} latent $\mathbf{x}$ is the object we probe.

\paragraph{Stage 2: latent-space diffusion world model.}
The world model is a next-frame autoregressive predictor. Given a window of
encoded context frames $\mathbf{x}_{0:N-1}$ and a noised (or fully masked)
target frame, it predicts the clean target latent $\mathbf{x}_N$. For the
continuous model this is done with a flow-matching objective realized by a
spatio-temporal diffusion transformer (STDiT): each block applies spatial
attention within a frame, temporal attention across frames, and a
feed-forward layer, while conditioning signals are injected through adaptive
layer normalization (AdaLN). At inference the model samples a noise tensor for
the target frame and integrates the learned velocity field back to a clean
latent; the generated frame is appended to the context and the earliest frame
is dropped, giving a sliding-window rollout. Optional ego-motion control
(steering angle and speed) is embedded by a small MLP and added to the AdaLN
conditioning, demonstrating that the latent dynamics can be steered.

Formally, the flow-matching objective defines a forward path from data to
noise by linear interpolation,
\begin{equation}
\mathbf{x}^{\tau} = (1-\tau)\,\mathbf{x} + \tau\,\boldsymbol{\epsilon},
\qquad \tau\in[0,1],\quad \boldsymbol{\epsilon}\sim\mathcal{N}(0,I),
\label{eq:fm-forward-intro}
\end{equation}
and trains a network to predict the velocity that transports the corrupted
target frame $\mathbf{x}_N^{\tau}$ back towards the data manifold,
conditioned on the context:
\begin{equation}
\mathcal{L} = \mathbb{E}_{\tau,\boldsymbol{\epsilon}}
\big[\,\lVert \mathbf{v}(\mathbf{x}_N^{\tau};\mathbf{x}_{0:N-1})
- (\boldsymbol{\epsilon}-\mathbf{x}_N) \rVert^2\,\big].
\label{eq:fm-loss-intro}
\end{equation}
We return to the theory behind Equations~\eqref{eq:fm-forward-intro}
and~\eqref{eq:fm-loss-intro} in Section~\ref{sec:diffusion-theory}.

\begin{figure}[t]
\centering
\begin{tikzpicture}[
  font=\small,
  box/.style={draw, rounded corners, minimum height=8mm, align=center, inner sep=3pt},
  enc/.style={box, fill=blue!8},
  lat/.style={box, fill=green!10, minimum width=9mm},
  wm/.style={box, fill=orange!12, minimum width=20mm, minimum height=15mm},
  dec/.style={box, fill=blue!8},
  ar/.style={-{Stealth[length=2mm]}, thick},
  >={Stealth[length=2mm]}
]
% Stage 1 label
\node[align=center, font=\small\bfseries] (s1) at (0,2.55) {Stage 1: Hybrid image tokenizer (VAE)};
% input image
\node[enc, fill=gray!12] (img) at (0,0) {Input\\frame $\mathbf{I}$};
\node[enc] (es) at (2.4,0.95) {Semantic\\enc.\ $\mathcal{E}_s$};
\node[enc] (ed) at (2.4,-0.95) {Detail\\enc.\ $\mathcal{E}_d$};
\node[lat] (xs) at (4.7,0.95) {$\mathbf{x}_s$};
\node[lat] (xd) at (4.7,-0.95) {$\mathbf{x}_d$};
\node[lat, minimum width=7mm, fill=green!18] (x) at (6.5,0) {$\mathbf{x}$};
\draw[ar] (img) -- (es);
\draw[ar] (img) -- (ed);
\draw[ar] (es) -- (xs);
\draw[ar] (ed) -- (xd);
\draw[ar] (xs) -- (x);
\draw[ar] (xd) -- (x);
\node[font=\scriptsize, align=center] at (2.4,1.95) {DINOv2\\cosine loss};
% Stage 2 label
\node[align=center, font=\small\bfseries] (s2) at (10.2,2.55) {Stage 2: Latent diffusion world model};
% context frames
\node[lat, fill=green!18, minimum width=6mm] (c0) at (8.2,0.9) {$\mathbf{x}_0$};
\node[font=\scriptsize] at (8.2,0.2) {$\cdots$};
\node[lat, fill=green!18, minimum width=6mm] (cN) at (8.2,-0.5) {$\mathbf{x}_{N\!-\!1}$};
\node[wm] (wm) at (10.7,0.2) {STDiT\\(flow\\matching)};
\node[lat, fill=green!18, minimum width=8mm] (xN) at (13.0,0.2) {$\hat{\mathbf{x}}_N$};
\draw[ar] (c0) -- (wm);
\draw[ar] (cN) -- (wm);
\draw[ar] (wm) -- (xN);
% noise into wm
\node[font=\scriptsize, align=center] (noise) at (10.4,-1.6) {noise $\boldsymbol{\epsilon}\sim\mathcal{N}(0,I)$\\+ AdaLN cond.};
\draw[ar] (noise) -- (wm);
% rollout feedback
\draw[ar, dashed] (xN.north) .. controls (13.0,2.0) and (8.2,2.0) .. (c0.north)
  node[midway, above, font=\scriptsize] {append, slide window};
% link stages: stage-1 latent feeds the stage-2 context window
\draw[ar, dashed] (x.east) .. controls (7.2,-0.1) and (7.4,-0.45) .. (cN.west);
% decoder
\node[dec] (dec) at (13.0,-1.8) {Image\\decoder};
\draw[ar] (xN) -- (dec);
\node[box, fill=gray!12] (out) at (13.0,-3.2) {Predicted\\next frame};
\draw[ar] (dec) -- (out);
\end{tikzpicture}
\caption{The two-stage Orbis architecture. \textbf{Stage~1} is a hybrid
tokenizer with a DINOv2-regularized semantic encoder and a detail encoder,
producing a continuous latent $\mathbf{x}$ (or, equivalently, discrete tokens).
\textbf{Stage~2} is a spatio-temporal diffusion transformer trained with flow
matching that predicts the next-frame latent $\hat{\mathbf{x}}_N$ from a window
of context latents; the prediction is appended and the window slides forward to
produce long rollouts. The frozen latent $\mathbf{x}$ is the representation we
probe in this work.}
\label{fig:orbis}
\end{figure}

\subsection{Probing the latent for future trajectories}
\label{sec:probe}

A standard way to interrogate what a learned representation ``knows'' is to
freeze it and train a small read-out head on a target task; if a simple probe
succeeds, the information is linearly or near-linearly accessible in the
representation~\cite{dosovitskiy_vit}. Our target task is \emph{future
trajectory prediction}: given the latent of the current (and recent) frames, we
want to predict the sequence of future 2D positions $(x,y)$ that the ego vehicle
will follow in its own coordinate frame. This is the same output space used by
modern end-to-end planners and by the trajectory-based evaluation of Orbis
itself, which makes the predictions directly comparable to driving benchmarks.

Two design choices define the probe. First, the read-out is an
\emph{attention} probe rather than a plain linear layer. The image latent is a
spatial grid of tokens, and which spatial regions matter for the future
trajectory --- lane markings, the vanishing point, leading vehicles --- depends
on the scene. A small set of learned query tokens attends over the frozen
latent grid and pools it into a fixed-size context vector, so the probe can
focus on the informative regions without modifying the backbone. Second, the
trajectory is produced by a \emph{diffusion} model operating directly in the 2D
trajectory space. Driving is inherently multi-modal: at an intersection,
turning left and turning right are both valid, and a regression head trained
with an $L_2$ loss would average them into an infeasible straight-ahead path.
A diffusion head instead learns the full conditional distribution of
trajectories and can sample several distinct, plausible futures, which is
exactly what is needed to assess the \emph{distribution} of behaviours encoded
in the world-model latent. Figure~\ref{fig:probe} summarizes the setup.

\begin{figure}[t]
\centering
\begin{tikzpicture}[
  font=\small,
  box/.style={draw, rounded corners, minimum height=8mm, align=center, inner sep=3pt},
  frozen/.style={box, fill=blue!8},
  train/.style={box, fill=orange!14},
  ar/.style={-{Stealth[length=2mm]}, thick},
  >={Stealth[length=2mm]}
]
% input
\node[box, fill=gray!12] (img) at (0,0) {Recent\\frames};
\node[frozen] (enc) at (2.3,0) {Frozen Orbis\\encoder $\mathcal{E}$};
% latent grid
\node[box, fill=green!12, minimum width=10mm, minimum height=12mm] (grid) at (4.7,0) {latent\\grid $\mathbf{x}$};
\draw[ar] (img) -- (enc);
\draw[ar] (enc) -- (grid);
\node[font=\scriptsize] at (2.3,-1.1) {(no gradient)};
% attention probe
\node[train, minimum height=12mm] (attn) at (7.4,0) {Attention\\probe\\(learned\\queries)};
\draw[ar] (grid) -- (attn) node[midway, above, font=\scriptsize] {keys/values};
\node[box, fill=green!18, minimum width=8mm] (ctx) at (9.8,0) {context\\$\mathbf{c}$};
\draw[ar] (attn) -- (ctx);
% diffusion head
\node[train, minimum height=12mm] (diff) at (12.0,0) {2D trajectory\\diffusion\\head};
\draw[ar] (ctx) -- (diff) node[midway, above, font=\scriptsize] {cond.};
\node[box, minimum height=12mm, fill=gray!10] (noise) at (12.0,-2.0) {noise\\$\to$ denoise};
\draw[ar] (noise) -- (diff);
% output trajectory
\begin{scope}[shift={(14.4,-0.1)}]
  \draw[->, gray] (-0.2,-0.7) -- (-0.2,0.9) node[above, font=\scriptsize, black] {$y$};
  \draw[->, gray] (-0.2,-0.7) -- (1.3,-0.7) node[right, font=\scriptsize, black] {$x$};
  \draw[blue, thick] (0.0,-0.6) .. controls (0.2,0.0) and (0.5,0.4) .. (0.9,0.7);
  \draw[red, thick, dashed] (0.0,-0.6) .. controls (0.0,0.1) and (-0.1,0.5) .. (0.0,0.8);
  \fill[blue] (0.0,-0.6) circle (1.2pt);
\end{scope}
\node[font=\scriptsize, align=center] at (15.0,-1.6) {sampled\\trajectories};
\end{tikzpicture}
\caption{The attentive diffusion probe. The Orbis encoder is frozen and
produces a spatial latent grid for the recent frames. A small attention probe
with learned query tokens pools the grid into a context vector $\mathbf{c}$,
which conditions a diffusion model that generates future 2D ego trajectories
$(x,y)$ directly in trajectory space. Because the head is generative, it can
sample several distinct futures (e.g.\ go straight vs.\ turn), letting us
evaluate the multi-modal distribution of trajectories rather than a single
averaged path.}
\label{fig:probe}
\end{figure}

This framing connects three threads that the rest of the report develops in
turn. The \emph{generative machinery} of both Orbis and the probe is diffusion
/ flow matching, reviewed in Section~\ref{sec:diffusion-theory}, and the
\emph{architecture} that carries it is the diffusion transformer of
Section~\ref{sec:dit}. The \emph{idea} of reusing a predictive video model as a
planning substrate places the work within the world-model literature of
Section~\ref{sec:world-models}, and in particular alongside navigation world
models that already use trajectory rollouts for decision making. Finally, the
\emph{evaluation} of the produced trajectories relies on the driving datasets
and planning metrics of Section~\ref{sec:datasets}.
